\documentclass[11pt]{article}
% Shared preamble for the printable mini lessons (1.4 onward).
%   \documentclass[11pt]{article}
%   % Shared preamble for the printable mini lessons (1.4 onward).
%   \documentclass[11pt]{article}
%   % Shared preamble for the printable mini lessons (1.4 onward).
%   \documentclass[11pt]{article}
%   \input{../preamble}
% Build the key with:  pdflatex -jobname=<name>_KEY "\def\iskey{}\input{<name>}"
% Figures shared with the web page live in figures/Unit_X/<name>.tex and are
% pulled in with \fig{Unit_X/<name>}; the standalone package drops their own
% preamble, so their macros come from figures/preamble.tex, loaded here.
\input{../../figures/preamble}
\usepackage[margin=0.75in]{geometry}
\usepackage{array}
\usepackage{enumitem}
\usepackage{amsmath}
\usepackage[helvet]{sfmath}  % math in the same sans font as the text
\usepackage{standalone}
\pagestyle{empty}
\setlength{\parindent}{0pt}
\setlist[enumerate,1]{leftmargin=1.5em,itemsep=1em}
\setlist[enumerate,2]{label=(\alph*),leftmargin=2em,itemsep=2.2em}

\newif\ifkey
\ifdefined\iskey \keytrue \fi
\colorlet{keyred}{red!75!black}
\newcommand{\ans}[1]{\ifkey\leavevmode{\color{keyred}#1}\fi}
% vertical work space on the student copy, less on the key
\newcommand{\work}[1]{\ifkey\vspace{0.6em}\else\vspace{#1}\fi}
\newcommand{\fig}[1]{\input{../../figures/#1}}
% Title line: \lessontitle{Quartiles, IQR, and Box Plots}
\newcommand{\lessontitle}[1]{{\Large\bfseries #1\ifkey\ \textcolor{keyred}{KEY}\fi}\hfill Name \underline{\hspace{2in}}\par\vspace{0.6em}}
\newcommand{\partner}{\par\vspace{0.4em}Work with a partner. Compare answers when you finish each problem.\par}

% Build the key with:  pdflatex -jobname=<name>_KEY "\def\iskey{}\input{<name>}"
% Figures shared with the web page live in figures/Unit_X/<name>.tex and are
% pulled in with \fig{Unit_X/<name>}; the standalone package drops their own
% preamble, so their macros come from figures/preamble.tex, loaded here.
% Shared preamble for web figures.  Each figure is a standalone .tex that
% \input's this file; scripts/make-figures.sh turns it into an SVG.
\usepackage[T1]{fontenc}
\usepackage{helvet}
\renewcommand{\familydefault}{\sfdefault}
\usepackage{tikz}
\usetikzlibrary{decorations.pathreplacing}
% axis numbers in the text font, not math italic
\pgfkeys{/pgf/number format/.cd,assume math mode=true,1000 sep={}}

% dot plot: \dotplot{min}{max}{value/count,...}{axis label}
\newcommand{\dotplot}[4]{%
\begin{tikzpicture}[x=0.9cm,y=0.33cm]
  \draw[thick] (#1-0.5,0) -- (#2+0.5,0);
  \foreach \t in {#1,...,#2} {\draw (\t,0.25) -- (\t,-0.25) node[below] {\t};}
  \foreach \v/\c in {#3} {\ifnum\c>0 \foreach \k in {1,...,\c} {\fill (\v,\k*1.15) circle (0.16cm);}\fi}
  \node at ({(#1+#2)/2},-2.3) {#4};
\end{tikzpicture}}

% shape sketch histogram (12 bins): \shapehist{heights}{label}
\newcommand{\shapehist}[2]{%
  \foreach \h [count=\i from 0] in {#1} \draw[fill=gray!30] (\i,0) rectangle (\i+1,\h);
  \draw (-0.3,0) -- (12.3,0);
  \node[font=\bfseries] at (6,-2.6) {#2};}
\newcommand{\meanmark}[1]{\fill (#1,0) -- ++(-0.3,-0.75) -- ++(0.6,0) -- cycle;}
\newcommand{\medmark}[2]{\draw[very thick,dashed] (#1,0) -- (#1,#2);}

% number line with minor ticks: \nline{min}{max}{label step}{axis label}
\newcommand{\nline}[4]{%
  \draw[thick] (#1-0.5,0) -- (#2+0.5,0);
  \foreach \t in {#1,...,#2} \draw (\t,0.12) -- (\t,-0.12);
  \pgfmathtruncatemacro{\nxt}{#1+#3}
  \foreach \t in {#1,\nxt,...,#2} \draw (\t,0.2) -- (\t,-0.2) node[below] {\t};
  \node at ({(#1+#2)/2},-0.95) {#4};}
% box plot above the line: \bplot{min}{Q1}{med}{Q3}{max}
\newcommand{\bplot}[5]{%
  \draw[thick] (#1,1.0) -- (#2,1.0);
  \draw[thick] (#4,1.0) -- (#5,1.0);
  \draw[thick] (#1,0.75) -- (#1,1.25);
  \draw[thick] (#5,0.75) -- (#5,1.25);
  \draw[thick,fill=gray!20] (#2,0.55) rectangle (#4,1.45);
  \draw[very thick] (#3,0.55) -- (#3,1.45);}

% ---- added for 1.4 onward ----
% axis on y=0: \axis{min}{max}{label step}{tick step}{axis label}
\newcommand{\axis}[5]{%
  \draw[thick] (#1,0) -- (#2,0);
  \pgfmathsetmacro{\axT}{#1+#4}\foreach \t in {#1,\axT,...,#2} \draw (\t,0.1) -- (\t,-0.1);
  \pgfmathsetmacro{\axL}{#1+#3}\foreach \t in {#1,\axL,...,#2}
    \draw (\t,0.18) -- (\t,-0.18) node[below] {\pgfmathprintnumber{\t}};
  \node at ({(#1+#2)/2},-0.95) {#5};}
% box plot centred at height y: \bplotat{y}{min}{Q1}{med}{Q3}{max}
\newcommand{\bplotat}[6]{%
  \draw[thick] (#2,#1) -- (#3,#1);
  \draw[thick] (#5,#1) -- (#6,#1);
  \draw[thick] (#2,#1-0.22) -- (#2,#1+0.22);
  \draw[thick] (#6,#1-0.22) -- (#6,#1+0.22);
  \draw[thick,fill=gray!20] (#3,#1-0.42) rectangle (#5,#1+0.42);
  \draw[very thick] (#4,#1-0.42) -- (#4,#1+0.42);}
% an outlier dot: \outlier{x}{y}
\newcommand{\outlier}[2]{\fill (#1,#2) circle (2.4pt);}

\usepackage[margin=0.75in]{geometry}
\usepackage{array}
\usepackage{enumitem}
\usepackage{amsmath}
\usepackage[helvet]{sfmath}  % math in the same sans font as the text
\usepackage{standalone}
\pagestyle{empty}
\setlength{\parindent}{0pt}
\setlist[enumerate,1]{leftmargin=1.5em,itemsep=1em}
\setlist[enumerate,2]{label=(\alph*),leftmargin=2em,itemsep=2.2em}

\newif\ifkey
\ifdefined\iskey \keytrue \fi
\colorlet{keyred}{red!75!black}
\newcommand{\ans}[1]{\ifkey\leavevmode{\color{keyred}#1}\fi}
% vertical work space on the student copy, less on the key
\newcommand{\work}[1]{\ifkey\vspace{0.6em}\else\vspace{#1}\fi}
\newcommand{\fig}[1]{\input{../../figures/#1}}
% Title line: \lessontitle{Quartiles, IQR, and Box Plots}
\newcommand{\lessontitle}[1]{{\Large\bfseries #1\ifkey\ \textcolor{keyred}{KEY}\fi}\hfill Name \underline{\hspace{2in}}\par\vspace{0.6em}}
\newcommand{\partner}{\par\vspace{0.4em}Work with a partner. Compare answers when you finish each problem.\par}

% Build the key with:  pdflatex -jobname=<name>_KEY "\def\iskey{}\input{<name>}"
% Figures shared with the web page live in figures/Unit_X/<name>.tex and are
% pulled in with \fig{Unit_X/<name>}; the standalone package drops their own
% preamble, so their macros come from figures/preamble.tex, loaded here.
% Shared preamble for web figures.  Each figure is a standalone .tex that
% \input's this file; scripts/make-figures.sh turns it into an SVG.
\usepackage[T1]{fontenc}
\usepackage{helvet}
\renewcommand{\familydefault}{\sfdefault}
\usepackage{tikz}
\usetikzlibrary{decorations.pathreplacing}
% axis numbers in the text font, not math italic
\pgfkeys{/pgf/number format/.cd,assume math mode=true,1000 sep={}}

% dot plot: \dotplot{min}{max}{value/count,...}{axis label}
\newcommand{\dotplot}[4]{%
\begin{tikzpicture}[x=0.9cm,y=0.33cm]
  \draw[thick] (#1-0.5,0) -- (#2+0.5,0);
  \foreach \t in {#1,...,#2} {\draw (\t,0.25) -- (\t,-0.25) node[below] {\t};}
  \foreach \v/\c in {#3} {\ifnum\c>0 \foreach \k in {1,...,\c} {\fill (\v,\k*1.15) circle (0.16cm);}\fi}
  \node at ({(#1+#2)/2},-2.3) {#4};
\end{tikzpicture}}

% shape sketch histogram (12 bins): \shapehist{heights}{label}
\newcommand{\shapehist}[2]{%
  \foreach \h [count=\i from 0] in {#1} \draw[fill=gray!30] (\i,0) rectangle (\i+1,\h);
  \draw (-0.3,0) -- (12.3,0);
  \node[font=\bfseries] at (6,-2.6) {#2};}
\newcommand{\meanmark}[1]{\fill (#1,0) -- ++(-0.3,-0.75) -- ++(0.6,0) -- cycle;}
\newcommand{\medmark}[2]{\draw[very thick,dashed] (#1,0) -- (#1,#2);}

% number line with minor ticks: \nline{min}{max}{label step}{axis label}
\newcommand{\nline}[4]{%
  \draw[thick] (#1-0.5,0) -- (#2+0.5,0);
  \foreach \t in {#1,...,#2} \draw (\t,0.12) -- (\t,-0.12);
  \pgfmathtruncatemacro{\nxt}{#1+#3}
  \foreach \t in {#1,\nxt,...,#2} \draw (\t,0.2) -- (\t,-0.2) node[below] {\t};
  \node at ({(#1+#2)/2},-0.95) {#4};}
% box plot above the line: \bplot{min}{Q1}{med}{Q3}{max}
\newcommand{\bplot}[5]{%
  \draw[thick] (#1,1.0) -- (#2,1.0);
  \draw[thick] (#4,1.0) -- (#5,1.0);
  \draw[thick] (#1,0.75) -- (#1,1.25);
  \draw[thick] (#5,0.75) -- (#5,1.25);
  \draw[thick,fill=gray!20] (#2,0.55) rectangle (#4,1.45);
  \draw[very thick] (#3,0.55) -- (#3,1.45);}

% ---- added for 1.4 onward ----
% axis on y=0: \axis{min}{max}{label step}{tick step}{axis label}
\newcommand{\axis}[5]{%
  \draw[thick] (#1,0) -- (#2,0);
  \pgfmathsetmacro{\axT}{#1+#4}\foreach \t in {#1,\axT,...,#2} \draw (\t,0.1) -- (\t,-0.1);
  \pgfmathsetmacro{\axL}{#1+#3}\foreach \t in {#1,\axL,...,#2}
    \draw (\t,0.18) -- (\t,-0.18) node[below] {\pgfmathprintnumber{\t}};
  \node at ({(#1+#2)/2},-0.95) {#5};}
% box plot centred at height y: \bplotat{y}{min}{Q1}{med}{Q3}{max}
\newcommand{\bplotat}[6]{%
  \draw[thick] (#2,#1) -- (#3,#1);
  \draw[thick] (#5,#1) -- (#6,#1);
  \draw[thick] (#2,#1-0.22) -- (#2,#1+0.22);
  \draw[thick] (#6,#1-0.22) -- (#6,#1+0.22);
  \draw[thick,fill=gray!20] (#3,#1-0.42) rectangle (#5,#1+0.42);
  \draw[very thick] (#4,#1-0.42) -- (#4,#1+0.42);}
% an outlier dot: \outlier{x}{y}
\newcommand{\outlier}[2]{\fill (#1,#2) circle (2.4pt);}

\usepackage[margin=0.75in]{geometry}
\usepackage{array}
\usepackage{enumitem}
\usepackage{amsmath}
\usepackage[helvet]{sfmath}  % math in the same sans font as the text
\usepackage{standalone}
\pagestyle{empty}
\setlength{\parindent}{0pt}
\setlist[enumerate,1]{leftmargin=1.5em,itemsep=1em}
\setlist[enumerate,2]{label=(\alph*),leftmargin=2em,itemsep=2.2em}

\newif\ifkey
\ifdefined\iskey \keytrue \fi
\colorlet{keyred}{red!75!black}
\newcommand{\ans}[1]{\ifkey\leavevmode{\color{keyred}#1}\fi}
% vertical work space on the student copy, less on the key
\newcommand{\work}[1]{\ifkey\vspace{0.6em}\else\vspace{#1}\fi}
\newcommand{\fig}[1]{\input{../../figures/#1}}
% Title line: \lessontitle{Quartiles, IQR, and Box Plots}
\newcommand{\lessontitle}[1]{{\Large\bfseries #1\ifkey\ \textcolor{keyred}{KEY}\fi}\hfill Name \underline{\hspace{2in}}\par\vspace{0.6em}}
\newcommand{\partner}{\par\vspace{0.4em}Work with a partner. Compare answers when you finish each problem.\par}


\begin{document}

\lessontitle{Venn Diagrams and the Addition Rule}

\textbf{Probability} is how likely something is, from 0 (impossible) to 1 (certain).
When every outcome is equally likely:
\[ P(A) = \frac{\text{number of ways } A \text{ can happen}}{\text{total number of outcomes}} \]
\vspace{-0.6em}%
$P(A)$ is read ``the probability of $A$.'' \qquad \textbf{Complement:} $P(\text{not } A) = 1 - P(A)$. \; (Some books write $P(A^C)$.)

\vspace{0.3em}
\textbf{Example:} A bag has 3 red and 7 blue marbles. Pick one. \; $P(\text{red}) = \dfrac{3}{10}$, and $P(\text{not red}) = 1 - \dfrac{3}{10} = \dfrac{7}{10}$.

\vspace{0.8em}
\textbf{Symbols for ``or'' and ``and''}

\vspace{0.2em}
\renewcommand{\arraystretch}{1.35}
\begin{tabular}{|c|l|l|}\hline
\textbf{Symbol} & \textbf{Read it as} & \textbf{Meaning}\\ \hline
$A \cup B$ & ``$A$ union $B$'' or ``$A$ or $B$'' & $A$ happens, or $B$ happens, or both\\ \hline
$A \cap B$ & ``$A$ intersect $B$'' or ``$A$ and $B$'' & $A$ and $B$ both happen (the overlap)\\ \hline
\end{tabular}

\vspace{0.8em}
\textbf{Example:} 30 students. 12 play soccer, 10 play basketball, 4 play both.

\begin{minipage}[c]{0.42\linewidth}\centering
\scalebox{0.72}{\fig{Unit_6/6.1_venn_sports}}
\end{minipage}\hfill
\begin{minipage}[c]{0.56\linewidth}
If we add $12 + 10$, the 4 students who play both are counted \textbf{twice}. Subtract them once.

\vspace{0.4em}
$P(S \cup B) = \dfrac{12}{30} + \dfrac{10}{30} - \dfrac{4}{30} = \dfrac{18}{30} = 0.6$
\end{minipage}

\vspace{0.6em}
\begin{tabular}{|>{\raggedright\arraybackslash}p{6.7in}|}\hline
\textbf{The addition rule:} \qquad $P(A \cup B) = P(A) + P(B) - P(A \cap B)$\\ \hline
\end{tabular}

\vspace{0.8em}
\textbf{Mutually exclusive} events cannot both happen, so $P(A \cap B) = 0$. Then $P(A \cup B) = P(A) + P(B)$.

\vspace{0.2em}
\textbf{Example:} Roll a number cube (1 to 6). Rolling a 1 and rolling a 2 are mutually exclusive.
\par $P(1 \cup 2) = \dfrac{1}{6} + \dfrac{1}{6} = \dfrac{2}{6} = \dfrac{1}{3}$

\vspace{0.8em}
\textbf{A deck of cards} has 52 cards: 4 suits (hearts, diamonds, clubs, spades), 13 cards in each suit (A, 2 to 10, J, Q, K).

\partner

% ---------------- PROBLEMS ----------------
\newpage
\begin{enumerate}

\item Roll a number cube (1 to 6). Find each probability.
\begin{enumerate}[itemsep=1.4em]
  \item $P(\text{even})$ \ans{\quad $\frac{3}{6} = \frac{1}{2}$}
  \item $P(\text{not } 6)$ \ans{\quad $1 - \frac{1}{6} = \frac{5}{6}$}
  \item $P(\text{greater than } 4)$ \ans{\quad $\frac{2}{6} = \frac{1}{3}$}
\end{enumerate}
\work{1em}

\item In a class of 40 students, 18 take Spanish ($S$), 15 take art ($A$), and 6 take both.
\begin{enumerate}[itemsep=1em]
  \item Fill in the Venn diagram.

  \begin{center}\scalebox{0.7}{\venn{Spanish ($S$)}{Art ($A$)}{\ans{12}}{\ans{6}}{\ans{9}}{\ans{13}}}\end{center}
  \item Find $P(S \cup A)$. \ans{\quad $\frac{18}{40} + \frac{15}{40} - \frac{6}{40} = \frac{27}{40}$}
  \work{1.5em}
  \item Find $P(S \cap A)$. \ans{\quad $\frac{6}{40} = \frac{3}{20}$}
  \item Find the probability that a student takes neither. \ans{\quad $\frac{13}{40}$}
\end{enumerate}
\work{1em}

\item $P(A) = 0.5$, \; $P(B) = 0.4$, and $P(A \cap B) = 0.2$. Find $P(A \cup B)$.

\ans{$0.5 + 0.4 - 0.2 = 0.7$}
\work{2.5em}

\newpage
\item Draw one card from a deck. Are the events mutually exclusive? Find the probability.
\begin{enumerate}[itemsep=1em]
  \item a heart or a spade

  \ans{Mutually exclusive (a card has one suit). \; $\frac{13}{52} + \frac{13}{52} = \frac{26}{52} = \frac{1}{2}$}
  \work{2.5em}
  \item a heart or a king

  \ans{Not mutually exclusive (the king of hearts). \; $\frac{13}{52} + \frac{4}{52} - \frac{1}{52} = \frac{16}{52} = \frac{4}{13}$}
\end{enumerate}
\work{2.5em}

\item Pick one student from this table at random.

\begin{center}
\renewcommand{\arraystretch}{1.4}
\begin{tabular}{|l|c|c|c|c|}\hline
 & \textbf{Bus} & \textbf{Car} & \textbf{Walk} & \textbf{Total} \\ \hline
\textbf{9th grade} & 18 & 12 & 10 & 40 \\ \hline
\textbf{10th grade} & 14 & 20 & 6 & 40 \\ \hline
\textbf{Total} & 32 & 32 & 16 & 80 \\ \hline
\end{tabular}
\end{center}
\begin{enumerate}[itemsep=1em]
  \item $P(\text{9th grade} \cap \text{car})$ \ans{\quad $\frac{12}{80} = 0.15$}
  \item $P(\text{10th grade} \cup \text{walk})$ \ans{\quad $\frac{40}{80} + \frac{16}{80} - \frac{6}{80} = \frac{50}{80} = 0.625$}
\end{enumerate}
\work{2em}

\item A student says, ``$P(A) = 0.6$ and $P(B) = 0.5$, so $P(A \cup B) = 1.1$.'' What did the student do wrong?

\ans{A probability cannot be more than 1. The student forgot to subtract $P(A \cap B)$, the overlap that was counted twice.}

\end{enumerate}

\end{document}
